\chapter*{แผนการสอนประจำวิชา}
\addcontentsline{toc}{chapter}{แผนบริหารการสอนประจำวิชา}
\vspace{0.3cm}
\noindent
\textbf{รหัสวิชา} \hspace{1cm}4091606 \\
\textbf{ชื่อวิชา} \hspace{1cm}คณิตศาสตร์สำหรับคอมพิวเตอร์ (Mathematics for Computer)\\
\textbf{จำนวนหน่วยกิต} \hspace{0.5cm}3(3-0-6)\\
\textbf{คณาจารย์ผู้สอน} \\
\indent ผู้ช่วยศาสตราจารย์ ดร.พิศิษฐ์ นาคใจ \quad หลักสูตรสาขาวิชาวิทยาการข้อมูล \quad ห้อง STA314 \\
\indent E-mail: pisit.nak@uru.ac.th \\[0.5em]
\textbf{คำอธิบายรายวิชา} \\
\indent
วิชา 4091606 คณิตศาสตร์สำหรับคอมพิวเตอร์ ศึกษาความรู้พื้นฐานเกี่ยวกับตรรกศาสตร์ เซต ความสัมพันธ์และฟังก์ชัน ระบบเลขฐานต่าง ๆ โดยเฉพาะเลขฐาน 2, 8 และ 16 เมทริกซ์และดีเทอร์มิแนนต์ ตลอดจนพีชคณิตบูลีน เพื่อใช้เป็นพื้นฐานสำหรับการคำนวณ การจัดการข้อมูล และการทำงานของระบบคอมพิวเตอร์

\begin{center}
\small
\begin{tabular}{|c|p{2.8cm}|p{4.5cm}|p{4.5cm}|}
\hline
\textbf{บท} & \textbf{ชื่อบท} & \textbf{รากฐานจากบทก่อน} & \textbf{สู่การประยุกต์ในบทถัดไป} \\ \hline
1 & ตรรกศาสตร์ (Logic) & --- & เป็นรากฐานให้ทุกบท \\ \hline
2 & ทฤษฎีเซต (Set Theory) & ตรรกศาสตร์ & เป็นรากฐานให้ทุกบท \\ \hline
3 & ความสัมพันธ์และฟังก์ชัน & บท 1--2 & ใช้ในบท 4, 5, 6 \\ \hline
4 & ระบบเลขฐาน & บท 1--2 & ใช้ในบท 6 \\ \hline
5 & เมทริกซ์และดีเทอร์มิแนนต์ & บท 1--2 & ใช้ในบท 6 \\ \hline
6 & พีชคณิตบูลีน & บท 1--5 ทั้งหมด & จุดสูงสุดของทุกบท \\ \hline
\end{tabular}
\end{center}

\textbf{บทที่ 1 ตรรกศาสตร์ (Logic)} เป็นบทแรกและเป็นรากฐานที่สำคัญที่สุดของวิชานี้ ศึกษาหลักการให้เหตุผลและการวิเคราะห์ค่าความจริงของประโยค โดยครอบคลุมประพจน์ ตัวเชื่อมประพจน์ (และ $\wedge$ หรือ $\vee$ นิเสธ $\neg$ ถ้า...แล้ว $\rightarrow$ ก็ต่อเมื่อ $\leftrightarrow$) ตารางค่าความจริง และการอนุมาน ตรรกศาสตร์เป็นเครื่องมือสำคัญในการออกแบบวงจรดิจิทัล การเขียนโปรแกรม และการพิสูจน์ความถูกต้องของขั้นตอนวิธี ในบริบทของการเชื่อมโยงกับบทอื่น หลักการตรรกศาสตร์จะถูกนำไปใช้ในบทที่ 3 สำหรับการวิเคราะห์ความสัมพันธ์ ในบทที่ 6 สำหรับการดำเนินการบูลีน (AND, OR, NOT) และในการพิสูจน์สูตรต่างๆ ตลอดทั้งวิชา

\textbf{บทที่ 2 ทฤษฎีเซต (Set Theory)} สร้างรากฐานต่อยอดจากตรรกศาสตร์โดยใช้แนวคิดเรื่องกลุ่มของสิ่งต่างๆ มาจัดระเบียบและจำแนกข้อมูล เนื้อหาครอบคลุมเซต สมาชิก การดำเนินการระหว่างเซต ได้แก่ ยูเนียน $A \cup B$ อินเตอร์เซกชัน $A \cap B$ ผลต่าง $A - B$ และคอมพลีเมนต์ $A^c$ รวมถึงเซตของจำนวน ทฤษฎีเซตเป็นภาษาพื้นฐานของฐานข้อมูลเชิงสัมพันธ์ คำสั่ง SQL ที่ใช้ UNION, INTERSECT, EXCEPT ล้วนสะท้อนการดำเนินการระหว่างเซตโดยตรง ในบทที่ 3 ความรู้เรื่องเซตจะถูกนำไปใช้กับผลคูณคาร์ทีเซียน $A \times B$ และความสัมพันธ์บนเซต สัญลักษณ์ $\cup$ $\cap$ $-$ $A^c$ จะปรากฏซ้ำในบทที่ 6 เมื่อเชื่อมโยงกับ AND, OR, NOT ในพีชคณิตบูลีน

\textbf{บทที่ 3 ความสัมพันธ์และฟังก์ชัน (Relations and Functions)} ต่อยอดจากทฤษฎีเซตโดยศึกษาความเชื่อมโยงระหว่างสมาชิกในเซตผ่านผลคูณคาร์ทีเซียน ความสัมพันธ์ $R \subseteq A \times B$ ช่วยให้เข้าใจโครงสร้างข้อมูลเชิงสัมพันธ์ในฐานข้อมูล กราฟในโครงสร้างข้อมูล และการวิเคราะห์เครือข่าย ฟังก์ชัน $f: A \rightarrow B$ ซึ่งเป็นความสัมพันธ์ชนิดพิเศษที่สมาชิกแต่ละตัวในเซตต้นกำเนิดจับคู่กับสมาชิกตัวเดียวในเซตเป้าหมาย มีบทบาทสำคัญในการเขียนโปรแกรมและการคำนวณ ในบริบทของบทที่ 5 ฟังก์ชันเชิงเส้นสามารถแทนด้วยเมทริกซ์ได้ และในบทที่ 6 ฟังก์ชันบูลีนเป็นพื้นฐานของลอจิกวงจรดิจิทัล

\textbf{บทที่ 4 ระบบเลขฐาน (Number Bases)} ศึกษาวิธีการแทนจำนวนในรูปแบบต่างๆ โดยเฉพาะฐาน 2 (binary) ฐาน 8 (octal) และฐาน 16 (hexadecimal) ซึ่งเป็นรากฐานของการทำงานของคอมพิวเตอร์ที่ใช้ระบบเลขฐานสองในการประมวลผลข้อมูล เนื้อหานี้เชื่อมโยงกับสถาปัตยกรรมคอมพิวเตอร์ การเข้ารหัสข้อมูล และการคำนวณทางคอมพิวเตอร์ในระดับต่ำ การแทนจำนวนในระบบเลขฐานต่างๆ ใช้หลักการเดียวกับการดำเนินการระหว่างเซตในบทที่ 2 คือการจัดกลุ่มและจำแนก ในบทที่ 6 ระบบเลขฐานสองเป็นพื้นฐานของการคำนวณบูลีนในวงจรดิจิทัล โดยบิต 0 และ 1 สอดคล้องกับค่าประพจน์เท็จและจริง

\textbf{บทที่ 5 เมทริกซ์และดีเทอร์มิแนนต์ (Matrices and Determinants)} ศึกษาโครงสร้างข้อมูลแบบตารางสองมิติพร้อมการดำเนินการทางคณิตศาสตร์ การบวก การคูณเมทริกซ์ และเมทริกซ์ผกผัน เมทริกซ์เป็นเครื่องมือสำคัญในคอมพิวเตอร์กราฟิกส์ การเรียนรู้ของเครื่อง และการประมวลผลภาพ การแปลงเมทริกซ์ใช้ในการหมุน การย่อ/ขยาย และการเลื่อนภาพ ความรู้เรื่องความสัมพันธ์จากบทที่ 3 ช่วยในการทำความเข้าใจการแปลงเชิงเส้นผ่านเมทริกซ์ และในบทที่ 6 เมทริกซ์จะถูกนำไปใช้ในการออกแบบวงจรคอมบิเนชันและซีเควนเชียล

\textbf{บทที่ 6 พีชคณิตบูลีน (Boolean Algebra)} เป็นบทสุดท้ายที่รวมความรู้จากทุกบทเข้าด้วยกัน การเชื่อมโยงระหว่างตรรกศาสตร์และทฤษฎีเซตปรากฏชัดเจนที่สุดในบทนี้ โดยค่าความจริง จริง (T) เท็จ (F) ในตรรกศาสตร์สอดคล้องกับสมาชิกในเซต และการดำเนินการบูลีน AND ($\cdot$) OR ($+$) NOT ($\bar{}$) มีคุณสมบัติเหมือนกับ $\cap$ $\cup$ complement ในทฤษฎีเซต กล่าวคือ $A \cap B$ มีคุณสมบัติเหมือน $A \cdot B$ และ $A \cup B$ มีคุณสมบัติเหมือน $A + B$ ตามกฎของเดมอร์แกน $\overline{A \cap B} = \bar{A} \cup \bar{B}$ ซึ่งมีรูปแบบเดียวกับในตรรกศาสตร์ พีชคณิตบูลีนเป็นรากฐานของวงจรดิจิทัล ตรรกศาสตร์ในการเขียนโปรแกรม และการค้นหาข้อมูลในฐานข้อมูลด้วยเงื่อนไข AND, OR, NOT

ทุกบทในวิชานี้เชื่อมโยงกันทั้ง 6 บท โดยบทที่ 1 และ 2 เป็นรากฐานสำคัญที่สุด หลักการจากทั้งสองบทนี้จะถูกนำไปใช้ในทุกบทถัดไป ความเข้าใจอย่างแท้จริงในตรรกศาสตร์และทฤษฎีเซตจะทำให้การเรียนรู้บทที่ 3 ถึง 6 เป็นไปอย่างราบรื่น และเมื่อมาถึงบทที่ 6 จะเห็นภาพรวมว่าทุกแนวคิดได้เชื่อมโยงกันอย่างเป็นระบบ
\\[0.5em]
\noindent
\textbf{จุดประสงค์การเรียนรู้}\\
\begin{enumerate}
    \item นักศึกษาสามารถอธิบายถึงความรู้พื้นฐานเกี่ยวกับตรรกศาสตร์ได้
    \item นักศึกษามีความรู้ความเข้าใจในคุณลักษณะของเซต พร้อมทั้งความสัมพันธ์และฟังก์ชันได้
    \item นักศึกษาสามารถแปลงและคำนวณในระบบเลขฐานต่าง ๆ โดยเฉพาะเลขฐาน 2, 8, 16 ได้
    \item นักศึกษามีความรู้ความเข้าใจในเมทริกซ์และดีเทอร์มิแนนต์ได้
    \item นักศึกษามีความรู้ความเข้าใจในพีชคณิตบูลีน ซึ่งเป็นเครื่องมือที่ทำให้เข้าใจการทำงานทางตรรกวิทยาภายในเครื่องคอมพิวเตอร์ได้
\end{enumerate}

%%%%
\noindent
\textbf{\normalsize โครงการสอน} \\
\indent
วิชาคณิตศาสตร์สำหรับคอมพิวเตอร์มีจำนวน 3 หน่วยกิต ใช้เวลาเรียน 3 ชั่วโมง/สัปดาห์ มีเวลาเรียนทั้งสิ้น 17 สัปดาห์ สามารถทำเป็นโครงการสอนได้ดังนี้

\newpage
\begingroup
\small
\setlength{\tabcolsep}{4pt}
\begin{longtable}{|c|p{4.3cm}|c|p{4.3cm}|c|}

\hline
\multicolumn{1}{|c|}{\textbf{สัปดาห์ที่}} &
\multicolumn{1}{c|}{\textbf{หัวข้อ/รายละเอียด}} &
\textbf{จำนวนชั่วโมง} &
\multicolumn{1}{c|}{\textbf{กิจกรรมการเรียนการสอน/สื่อที่ใช้}} &
\multicolumn{1}{c|}{\textbf{ผู้สอน}} \\ \hline
\endfirsthead
\hline
\multicolumn{1}{|c|}{\textbf{สัปดาห์ที่}} &
\multicolumn{1}{c|}{\textbf{หัวข้อ/รายละเอียด}} &
\textbf{จำนวนชั่วโมง} &
\multicolumn{1}{c|}{\textbf{กิจกรรมการเรียนการสอน/สื่อที่ใช้}} &
\multicolumn{1}{c|}{\textbf{ผู้สอน}} \\ \hline
\endhead

1--2 &
บทที่ 1 ความรู้พื้นฐานเกี่ยวกับตรรกศาสตร์ &
6 &
\begin{tabular}[t]{@{}l@{}}บรรยาย สรุปบทเรียน\\ อภิปราย ซักถาม และ\\ แลกเปลี่ยนเรียนรู้\\ การตอบคำถามท้ายบท\end{tabular} &
\begin{tabular}[t]{@{}l@{}}ผศ.ดร.พิศิษฐ์\\ นาคใจ\end{tabular} \\ \hline

3--4 &
บทที่ 2 เซต &
6 &
\begin{tabular}[t]{@{}l@{}}บรรยาย สรุปบทเรียน\\ อภิปราย ซักถาม และ\\ แลกเปลี่ยนเรียนรู้\\ ศึกษาเอกสารคำสอน\\ การตอบคำถามท้ายบท\end{tabular} &
\begin{tabular}[t]{@{}l@{}}ผศ.ดร.พิศิษฐ์\\ นาคใจ\end{tabular} \\ \hline

5--7 &
บทที่ 3 ความสัมพันธ์และฟังก์ชัน &
9 &
\begin{tabular}[t]{@{}l@{}}บรรยาย สรุปบทเรียน\\ อภิปราย ซักถาม และ\\ แลกเปลี่ยนเรียนรู้\\ ศึกษาเอกสารคำสอน\\ การตอบคำถามท้ายบท\end{tabular} &
\begin{tabular}[t]{@{}l@{}}ผศ.ดร.พิศิษฐ์\\ นาคใจ\end{tabular} \\ \hline

8 & \multicolumn{4}{c|}{\textbf{สอบกลางภาคเรียน}} \\ \hline

9--11 &
บทที่ 4 ระบบเลขฐานต่าง ๆ โดยเฉพาะเลขฐาน 2, 8, 16 &
9 &
\begin{tabular}[t]{@{}l@{}}บรรยาย สรุปบทเรียน\\ อภิปราย ซักถาม และ\\ แลกเปลี่ยนเรียนรู้\\ ศึกษาเอกสารคำสอน\\ การตอบคำถามท้ายบท\end{tabular} &
\begin{tabular}[t]{@{}l@{}}ผศ.ดร.พิศิษฐ์\\ นาคใจ\end{tabular} \\ \hline

12--13 &
บทที่ 5 เมทริกซ์และดีเทอร์มิแนนต์ &
6 &
\begin{tabular}[t]{@{}l@{}}บรรยาย สรุปบทเรียน\\ อภิปราย ซักถาม และ\\ แลกเปลี่ยนเรียนรู้\\ ศึกษาเอกสารคำสอน\\ การตอบคำถามท้ายบท\end{tabular} &
\begin{tabular}[t]{@{}l@{}}ผศ.ดร.พิศิษฐ์\\ นาคใจ\end{tabular} \\ \hline

14--16 &
บทที่ 6 พีชคณิตบูลีน &
9 &
\begin{tabular}[t]{@{}l@{}}บรรยาย สรุปบทเรียน\\ อภิปราย ซักถาม และ\\ แลกเปลี่ยนเรียนรู้\\ ศึกษาเอกสารคำสอน\\ การตอบคำถามท้ายบท\end{tabular} &
\begin{tabular}[t]{@{}l@{}}ผศ.ดร.พิศิษฐ์\\ นาคใจ\end{tabular} \\ \hline

17 & \multicolumn{4}{c|}{\textbf{สอบปลายภาคเรียน}} \\ \hline

\end{longtable}
\endgroup
\addtocounter{table}{-1}%<-- decrement table counter

\newpage
\noindent
\textbf{\normalsize วิธีสอนและกิจกรรมการเรียนรู้}
\begin{enumerate}
    \item วิธีสอน
    \begin{enumerate}
        \item วิธีสอนส่วนใหญ่ใช้วิธีการบรรยายประกอบเอกสารคำสอน และจัดกิจกรรมให้นักศึกษาอภิปราย ซักถาม และแลกเปลี่ยนเรียนรู้ในชั้นเรียน
    \end{enumerate}
    \item กิจกรรมการเรียนรู้
    \begin{enumerate}
        \item ให้นักศึกษาตอบคำถามประจำบทในเอกสารคำสอน
        \item ให้นักศึกษาศึกษาเอกสารคำสอนเพิ่มเติมตามหัวข้อที่กำหนด
        \item ให้นักศึกษาร่วมอภิปราย ซักถาม และแลกเปลี่ยนเรียนรู้ในชั้นเรียน
    \end{enumerate}
\end{enumerate}

\noindent
\textbf{\normalsize สื่อการเรียนรู้}
\begin{enumerate}
    \item เอกสารคำสอน ตำรา และบทความวิจัยที่เกี่ยวข้อง
    \item สื่ออิเล็กทรอนิกส์ เช่น เว็บไซต์และแหล่งเรียนรู้ออนไลน์ต่าง ๆ
    \item คำถามประจำบท
\end{enumerate}

%%%%%%%%%%%%%%%%%%%%%%%
%การวัดและการประเมินผล
\noindent
\textbf{\normalsize การวัดและการประเมินผล}
\begin{enumerate}
    \item คะแนนระหว่างภาคเรียน \hfill 70\%
    \begin{enumerate}
        \item การตรงต่อเวลาและความซื่อสัตย์ \hfill 10\%
        \item งานที่ได้รับมอบหมาย (Assignment) \hfill 10\%
        \item กิจกรรมกลุ่มและการนำเสนอ \hfill 20\%
        \item ทดสอบกลางภาคเรียน \hfill 30\%
    \end{enumerate}
    \item ทดสอบปลายภาคเรียน \hfill 30\%
\end{enumerate}

\begin{table}[h]
    \centering
    \begin{tabular}{ccc}
        \hline
        \textbf{คะแนน} & \textbf{เกรด} & \textbf{ความหมาย} \\
        \hline
        80--100  & A   & ดีเยี่ยม   \\
        75--79   & B$^{+}$ & ดีมาก  \\
        70--74   & B   & ดี         \\
        65--69   & C$^{+}$ & ดีพอใช้ \\
        60--64   & C   & พอใช้      \\
        55--59   & D$^{+}$ & อ่อน    \\
        50--54   & D   & อ่อนมาก   \\
        0--49    & F   & ตก         \\
        \hline
    \end{tabular}
\end{table}

\noindent
\textbf{\normalsize ตำรา เอกสาร และงานวิจัยประกอบการเรียนการสอน}
\begin{enumerate}
    \item เอกสารหลัก \\
    คณะกรรมการโครงการตำราวิทยาศาสตร์และคณิตศาสตร์มูลนิธิ สอวน. (2556).
    \textit{คณิตศาสตร์พื้นฐานสำหรับคอมพิวเตอร์}. พิมพ์ครั้งที่ 4. กรุงเทพฯ: บริษัท ด่านสุทธาการพิมพ์ จำกัด.

    \item เอกสารเพิ่มเติม
    \begin{enumerate}
        \item คณะกรรมการโครงการตำราวิทยาศาสตร์และคณิตศาสตร์มูลนิธิ สอวน. (2553).
        \textit{คณิตศาสตร์พื้นฐานสำหรับคอมพิวเตอร์}. พิมพ์ครั้งที่ 3. กรุงเทพฯ: บริษัท ด่านสุทธาการพิมพ์ จำกัด.
        \item อุ่นใจ ลิมตระกูล. (2538). \textit{คณิตศาสตร์คอมพิวเตอร์}. พิมพ์ครั้งที่ 1. กรุงเทพฯ: สำนักพิมพ์ภูมิบัณฑิต.
    \end{enumerate}

    \item งานวิจัยของผู้สอนที่ใช้ประกอบการเรียนการสอน \\
    งานวิจัยต่อไปนี้ใช้เนื้อหาในรายวิชาเป็นเครื่องมือหลัก จึงนำมาใช้เป็นกรณีศึกษาประกอบการเรียนการสอนตามความเหมาะสม
    \begin{enumerate}
        \item Nakjai, P., \& Katanyukul, T. (2019). Hand sign recognition for Thai finger spelling:
        An application of convolution neural network. \textit{Journal of Signal Processing Systems},
        91(2), 131--146. \url{https://doi.org/10.1007/s11265-018-1375-6}

        \item Nakjai, P., Maneerat, P., \& Katanyukul, T. (2019). Thai finger spelling localization
        and classification under complex background using a YOLO-based deep learning.
        \textit{Proceedings of the 11th International Conference on Computer Modeling and Simulation}.

        \item Nakjai, P., \& Katanyukul, T. (2021). Automatic Thai finger spelling transcription.
        \textit{Walailak Journal of Science and Technology}, 18(10).
        \url{https://doi.org/10.48048/wjst.2021.11233}

        \item Nakjai, P., \& Katanyukul, T. (2018). Automatic hand sign recognition: Identify
        unusuality through latent cognizance. \textit{Artificial Neural Networks in Pattern
        Recognition (ANNPR 2018)}, Lecture Notes in Computer Science, 11081.
        \url{https://doi.org/10.1007/978-3-319-99978-4_20}

        \item Nakjai, P., \& Katanyukul, T. (2022). Toward latent cognizance on open-set recognition.
        \textit{Intelligent Information and Database Systems}, Lecture Notes in Computer Science.
        \url{https://doi.org/10.1007/978-3-030-98018-4_20}

        \item Katanyukul, T., \& Nakjai, P. (2022). Recognition awareness: Adding awareness to
        pattern recognition using latent cognizance. \textit{Heliyon}, 8(4), e09240.
        \url{https://doi.org/10.1016/j.heliyon.2022.e09240}
    \end{enumerate}
\end{enumerate}

\vspace{2em}
\noindent
\begin{tabular}{ll}
ลงชื่อ & .......ผู้ช่วยศาสตราจารย์ ดร.พิศิษฐ์ นาคใจ..... อาจารย์ผู้สอน \\
\end{tabular}
